\documentclass[11pt,a4paper]{article}
\usepackage[utf8]{inputenc}
\usepackage[spanish,es-nodecimaldot]{babel}
\usepackage[margin=2.2cm]{geometry}
\usepackage{amsmath,amssymb,amsfonts}
\usepackage{booktabs}
\usepackage{xcolor}
\usepackage{hyperref}
\usepackage{fancyhdr}
\usepackage{tcolorbox}
\usepackage{titlesec}
\usepackage{enumitem}

\definecolor{primary}{RGB}{15, 23, 42}
\definecolor{accent}{RGB}{37, 99, 235}
\definecolor{boxbg}{RGB}{248, 250, 252}
\definecolor{boxborder}{RGB}{226, 232, 240}
\definecolor{headerbg}{RGB}{30, 58, 138}

\hypersetup{
    colorlinks=true,
    linkcolor=accent,
    filecolor=accent,
    urlcolor=accent,
    pdftitle={Memoria Técnica y Formulación Matemática - SASID CDMX},
    pdfauthor={SASID Pro - Sobrio25}
}

\pagestyle{fancy}
\fancyhf{}
\fancyhead[L]{\small\textsf{\textcolor{gray}{SASID Pro | Memoria Técnica y Formulación Matemática}}}
\fancyhead[R]{\small\textsf{\textcolor{gray}{NTC-Sismo CDMX 2017/2023}}}
\fancyfoot[C]{\thepage}
\renewcommand{\headrulewidth}{0.4pt}

\titleformat{\section}{\Large\bfseries\color{primary}}{\thesection.}{0.6em}{}[\titlerule]
\titleformat{\subsection}{\large\bfseries\color{headerbg}}{\thesubsection.}{0.5em}{}

\tcbuselibrary{skins,breakable}
\newtcolorbox{formulabox}[1]{
    colback=boxbg,
    colframe=accent,
    fonttitle=\bfseries\color{white},
    title=#1,
    boxrule=1pt,
    arc=4pt,
    left=8pt,
    right=8pt,
    top=6pt,
    bottom=6pt,
    breakable
}

\begin{document}

% Encabezado Principal
\begin{center}
    {\Huge\bfseries\color{primary} Formulación Matemática y Memoria Técnica}\\[6pt]
    {\LARGE\bfseries\color{headerbg} Generación de Espectros Sísmicos de Diseño (NTC-CDMX)}\\[10pt]
    {\large\textbf{Sistema de Acciones Sísmicas de Diseño -- SASID Pro}}\\[4pt]
    \textsf{\textcolor{gray}{Implementación en Flutter: \href{https://github.com/Sobrio25/sasid-pro}{github.com/Sobrio25/sasid-pro}}}\\[14pt]
    \rule{\linewidth}{1pt}
\end{center}

\vspace{0.2cm}

\section{Introducción y Marco Normativo}
El presente documento describe de manera rigurosa las expresiones matemáticas, los parámetros geotécnicos y los factores estructurales implementados en la aplicación \textbf{SASID Pro} para la construcción de los espectros de respuesta elásticos e inelásticos de diseño sísmico en la Ciudad de México.

La metodología y formulación analítica se sustentan en los siguientes instrumentos oficiales:
\begin{itemize}[noitemsep,topsep=2pt]
    \item \textbf{Normas Técnicas Complementarias para Diseño por Sismo (NTC-Sismo 2017 / 2023)}, publicadas en la \textit{Gaceta Oficial de la Ciudad de México} (Apéndice A y Capítulo 3).
    \item \textbf{Manual Técnico SASID A}, desarrollado por el Instituto de Ingeniería de la UNAM (\textsf{II-UNAM}) y la Sociedad Mexicana de Ingeniería Sísmica (\textsf{SMIS}).
    \item \textbf{NTC-Sismo 2004}, aplicadas en el módulo de evaluación comparativa histórica.
\end{itemize}

\vspace{0.2cm}

\section{Zonificación Geotécnica y Parámetros del Sitio}
La demanda sísmica en la Cuenca de México depende fuertemente de las condiciones estratigráficas locales, caracterizadas por el periodo dominante de vibración del suelo $T_s$. Las NTC clasifican el terreno en las siguientes zonas:

\begin{table}[h!]
\centering
\small
\renewcommand{\arraystretch}{1.25}
\begin{tabular}{@{}llcccl@{}}
\toprule
\textbf{Zona} & \textbf{Denominación} & \textbf{Rango $T_s$ (s)} & \textbf{$a_0$ (g)} & \textbf{$c$ (g)} & \textbf{Descripción Geotécnica} \\
\midrule
\textbf{Zona I}   & Lomas (Firme)         & $T_s \le 0.5$    & $0.06$       & $0.18 - 0.22$ & Tobas volcánicas, arenas y gravas compactas. \\
\textbf{Zona II}  & Transición            & $0.5 < T_s \le 1.0$ & $0.08 - 0.11$ & $0.22 - 0.45$ & Depósitos arcillosos de espesor moderado. \\
\textbf{Zona IIIa}& Lago (Blando)         & $1.0 < T_s \le 1.5$ & $0.11 - 0.14$ & $0.45 - 0.65$ & Arcillas altamente compresibles. \\
\textbf{Zona IIIb}& Lago (Espesor medio)  & $1.5 < T_s \le 2.0$ & $0.14 - 0.16$ & $0.65 - 0.85$ & Arcillas lacustres de gran espesor. \\
\textbf{Zona IIIc}& Lago (Profundo)       & $2.0 < T_s \le 3.0$ & $0.16 - 0.19$ & $0.85 - 1.00$ & Depósitos arcillosos profundos. \\
\textbf{Zona IIId}& Lago (Muy blando)     & $T_s > 3.0$      & $0.19 - 0.22$ & $1.00 - 1.10$ & Zonas de máxima amplificación dinámica. \\
\bottomrule
\end{tabular}
\caption{Parámetros de sitio según la zonificación geotécnica de las NTC-CDMX.}
\end{table}

\vspace{-0.2cm}

\section{Formulación del Espectro Elástico Transparente $a_e(T)$}
El espectro elástico representa la seudoaceleración máxima normalizada ($S_a / g$) que experimenta un oscilador simple con una fracción de amortiguamiento crítico $\zeta = 5\%$. De acuerdo con el \textbf{Apéndice A} de las NTC-Sismo:

\begin{formulabox}{Ecuación del Espectro Elástico $a_e(T)$ (Apéndice A NTC)}
$$
a_e(T) = 
\begin{cases} 
a_0 + (c - a_0) \left(\dfrac{T}{T_a}\right) & \text{si } T < T_a \quad \text{(Rama inicial ascendente)} \\[10pt]
c & \text{si } T_a \le T \le T_b \quad \text{(Meseta espectral)} \\[10pt]
c \left(\dfrac{T_b}{T}\right)^k & \text{si } T > T_b \quad \text{(Rama descendente de periodos largos)}
\end{cases}
$$
\end{formulabox}

\textbf{Nomenclatura de variables:}
\begin{itemize}[noitemsep]
    \item $T$: Periodo fundamental de vibración de la estructura ($0.0 \le T \le 5.0\,\text{s}$).
    \item $a_0$: Aceleración máxima del terreno (\textit{Peak Ground Acceleration}, PGA) en fracción de la gravedad $g$.
    \item $c$: Coeficiente sísmico elástico (ordenada máxima constante de la meseta).
    \item $T_a, T_b$: Periodos característicos que delimitan el inicio y fin de la meseta espectral.
    \item $k$: Exponente de decaimiento espectral que controla la caída de aceleraciones para periodos largos.
\end{itemize}

\newpage

\section{Factores de Modificación y Reducción Inelástica}
Para considerar la capacidad de disipación de energía por ductilidad y sobrerresistencia de la estructura, las ordenadas elásticas se reducen mediante los factores reglamentarios:

\subsection{1. Factor de Ductilidad Dependiente del Periodo $Q'(T)$}
El factor de comportamiento sísmico nominal $Q$ ($1.0, 1.5, 2.0, 3.0, 4.0$) se corrige en función del periodo estructural para evitar reducciones excesivas en estructuras muy rígidas ($T < T_a$):

\begin{formulabox}{Factor de Comportamiento Sísmico Reducido $Q'(T)$}
$$
Q'(T) = 
\begin{cases} 
1 + (Q - 1) \left(\dfrac{T}{T_a}\right) & \text{si } T < T_a \\[8pt]
Q & \text{si } T \ge T_a
\end{cases}
$$
\end{formulabox}

\subsection{2. Factor de Sobrerresistencia Dependiente del Periodo $R(T)$}
La sobrerresistencia estructural considera la reserva plástica y la redistribución de esfuerzos:

\begin{formulabox}{Factor de Sobrerresistencia Reducido $R(T)$}
$$
R(T) = 
\begin{cases} 
1 + (R_0 - 1) \left(\dfrac{T}{T_a}\right) & \text{si } T < T_a \\[8pt]
R_0 & \text{si } T \ge T_a
\end{cases}
$$
\end{formulabox}

Donde el factor de sobrerresistencia base $R_0$ se calcula a partir de $Q$ y del factor de hiperestaticidad $k_1$:
\begin{equation}
R_0 = 
\begin{cases}
2.0 \cdot k_1 & \text{para } Q \ge 3.0 \\
1.75 \cdot k_1 & \text{para } Q = 2.0 \\
1.50 \cdot k_1 & \text{para } Q = 1.5 \\
1.00 & \text{para } Q = 1.0
\end{cases}
\end{equation}

\subsection{3. Factor de Importancia ($I$) y Factor de Irregularidad ($\alpha$)}
\begin{itemize}[noitemsep]
    \item \textbf{Factor de Importancia ($I$):}
    \begin{itemize}[noitemsep]
        \item \textbf{Grupo A} ($I = 1.50$): Hospitales, escuelas, centrales eléctricas, estaciones de bomberos.
        \item \textbf{Grupo B} ($I = 1.00$): Viviendas, comercios, oficinas, hoteles e industrias estándar.
        \item \textbf{Grupo A1} ($I = 1.75$ / NTC 2023): Instalaciones críticas con materiales peligrosos.
    \end{itemize}
    \item \textbf{Factor de Irregularidad ($\alpha$):}
    \begin{itemize}[noitemsep]
        \item Regular: $\alpha = 1.00$ \quad | \quad Irregular: $\alpha = 0.90$
        \item Fuertemente Irregular: $\alpha = 0.80$ \quad | \quad Extremadamente Irregular: $\alpha = 0.70$
    \end{itemize}
\end{itemize}

\vspace{0.2cm}

\section{Formulación del Espectro Inelástico de Diseño $a_d(T)$}
El espectro de diseño final aplicable al análisis modal espectral se calcula afectando la aceleración elástica por los factores anteriores, sujeto a una cota inferior obligatoria:

\begin{formulabox}{Espectro Inelástico de Diseño $a_d(T)$}
$$
a_d(T) = \max \left( \frac{I \cdot a_e(T)}{Q'(T) \cdot R(T) \cdot \alpha}, \; a_{\text{mín}} \right)
$$
\end{formulabox}

Donde la aceleración mínima de diseño reglamentaria $a_{\text{mín}}$ se evalúa como:
\begin{equation}
a_{\text{mín}} = \frac{a_0 \cdot I}{2 \cdot R_0 \cdot \alpha}
\end{equation}

\vspace{0.2cm}

\section{Espectro de Peligro Uniforme (EPU)}
El Espectro de Peligro Uniforme corresponde a la envolvente probabilista con periodo de retorno de diseño (475 años) sin afectación por factores inelásticos:
\begin{equation}
a_{\text{epu}}(T) = 1.10 \cdot I \cdot a_e(T)
\end{equation}

\newpage

\section{Formulación Comparativa: NTC-Sismo 2004}
Para efectos de contraste histórico en la toma de decisiones estructurales, el programa incluye las formulaciones de las normas de 2004:

\subsection{1. Cuerpo Principal NTC 2004 (Por Zonas Clásicas)}
$$
a_{2004}(T) = 
\begin{cases} 
a_0 + (c - a_0) \left(\dfrac{T}{T_a}\right) & \text{si } T < T_a \\[6pt]
c & \text{si } T_a \le T \le T_b \\[6pt]
c \left(\dfrac{T_b}{T}\right)^r & \text{si } T > T_b 
\end{cases}
\qquad \implies \qquad
a_{d,2004}(T) = \frac{a_{2004}(T)}{Q'(T)}
$$

\begin{table}[h!]
\centering
\small
\renewcommand{\arraystretch}{1.15}
\begin{tabular}{lccccc}
\toprule
\textbf{Zona NTC 2004} & \textbf{$c$} & \textbf{$a_0$} & \textbf{$T_a$ (s)} & \textbf{$T_b$ (s)} & \textbf{Exponente $r$} \\
\midrule
Zona I (Lomas)        & 0.16 & 0.04 & 0.20 & 0.60 & 0.50 \\
Zona II (Transición)  & 0.32 & 0.08 & 0.30 & 1.50 & 0.67 \\
Zona III (Lago)       & 0.45 & 0.10 & 0.60 & 2.50 & 1.00 \\
\bottomrule
\end{tabular}
\caption{Parámetros normativos del cuerpo principal de las NTC-Sismo 2004.}
\end{table}

\vspace{0.2cm}

\section{Correspondencia en el Repositorio de Código}
La siguiente tabla indica la ubicación exacta de cada componente matemático en el repositorio \href{https://github.com/Sobrio25/sasid-pro}{\texttt{Sobrio25/sasid-pro}}:

\begin{table}[h!]
\centering
\small
\renewcommand{\arraystretch}{1.3}
\begin{tabular}{@{}lll@{}}
\toprule
\textbf{Fórmula / Módulo} & \textbf{Archivo en el Repositorio} & \textbf{Líneas / Función} \\
\midrule
Interpolación Geotécnica ($T_s, a_0, c, T_a, T_b, k$) & \texttt{lib/services/seismic\_engine.dart} & \texttt{calculateSiteParameters()} (L7--114) \\
Espectro Elástico $a_e(T)$ & \texttt{lib/services/seismic\_engine.dart} & \texttt{computeSpectrum()} (L145--153) \\
Factores Inelásticos $Q'(T), R(T), R_0$ & \texttt{lib/services/seismic\_engine.dart} & \texttt{computeSpectrum()} (L118--165) \\
Espectro de Diseño $a_d(T)$ y $a_{\text{mín}}$ & \texttt{lib/services/seismic\_engine.dart} & \texttt{computeSpectrum()} (L167--177) \\
Espectro de Peligro Uniforme (EPU) & \texttt{lib/services/seismic\_engine.dart} & \texttt{computeSpectrum()} (L180--182) \\
Comparativa Multicriterio ($Q=1.5, 2, 3, 4$) & \texttt{lib/services/seismic\_engine.dart} & \texttt{\_computeDesignForQ()} (L197--233) \\
Espectros NTC-2004 (Cuerpo y Apéndice A) & \texttt{lib/services/seismic\_engine.dart} & \texttt{\_computeNtc2004Main()} (L235--302) \\
Modelos de Datos y Constantes Normativas & \texttt{lib/models/seismic\_models.dart} & Definiciones de Enums y Clases \\
Exportadores TXT (SAP2000), CSV y PDF & \texttt{lib/services/export\_service.dart} & \texttt{ExportService} (L8--180) \\
Pruebas Unitarias de Validación Numérica & \texttt{test/seismic\_engine\_test.dart} & Suite de pruebas de precisión \\
\bottomrule
\end{tabular}
\caption{Mapeo de formulaciones matemáticas contra el código fuente en Dart.}
\end{table}

\vspace{0.4cm}
\begin{center}
    \rule{0.6\linewidth}{0.5pt}\\[6pt]
    \small\textsf{\textbf{SASID Pro} -- Desarrollado en Flutter para Ingeniería Sísmica y Estructural}\\
    \small\textsf{Ciudad de México $\cdot$ 2026}
\end{center}

\end{document}
